% -------------------------
% Page layout and typography
% -------------------------
\usepackage[a4paper,inner=26mm,outer=26mm,top=30mm,bottom=32mm,headheight=15pt]{geometry}
\usepackage{setspace}
\onehalfspacing

\usepackage{iftex}
\usepackage{soul}

\ifPDFTeX
  \usepackage[T1]{fontenc}
  \usepackage[utf8]{inputenc}
  \usepackage{newtxtext,newtxmath}
\else
  \usepackage{fontspec}
  \setmainfont{Libertinus Serif}
  \setsansfont{Libertinus Sans}
  \setmonofont{Libertinus Mono}
\fi

\usepackage{microtype}
\usepackage{csquotes}

% -------------------------
\usepackage{tcolorbox}
\newtcolorbox{questionbox}{
    colback=RedViolet!5!white,
    colframe=RedViolet!75!black,
    boxrule=0.8pt,
    arc=0pt,
    left=6pt,
    right=6pt,
    top=6pt,
    bottom=6pt
}

% resultbox
\usepackage[dvipsnames]{xcolor}
\usepackage{tcolorbox}

\newtcolorbox{resultbox}{
    colback=YellowGreen!5!white,
    colframe=YellowGreen!75!black,
    boxrule=0.8pt,
    arc=0pt,
    left=6pt,
    right=6pt,
    top=6pt,
    bottom=6pt
}

%%% 

% \usepackage[a4paper,margin=1in]{geometry}
% \usepackage{amsmath,amssymb,amsthm,mathtools}
% \usepackage{booktabs}
% \usepackage{enumitem}
% \usepackage{hyperref}
% \usepackage{xcolor}
% \usepackage[most]{tcolorbox}


% -------------------------
% Mathematics
% -------------------------
\usepackage{amsmath,amssymb,amsfonts,mathtools,bm,bbm}
% \usepackage{amssymb}
\usepackage{commath}
\usepackage{physics}
\usepackage{siunitx}
\usepackage{enumitem}
\numberwithin{equation}{chapter}

% Common operators/macros for stochastic analysis and finance
\DeclareMathOperator{\Tr}{Tr}
\DeclareMathOperator{\diag}{diag}
\DeclareMathOperator{\Cov}{Cov}
\DeclareMathOperator{\Var}{Var}
\DeclareMathOperator{\Corr}{Corr}
\DeclareMathOperator{\esssup}{ess\,sup}
\newcommand{\R}{\mathbb{R}}
\newcommand{\N}{\mathbb{N}}
\newcommand{\E}{\mathbb{E}}
\newcommand{\Prob}{\mathbb{P}}
\newcommand{\Q}{\mathbb{Q}}
\newcommand{\F}{\mathcal{F}}
\newcommand{\dd}{\,\mathrm{d}}
\newcommand{\ii}{\mathrm{i}}
\newcommand{\SPD}{\mathbb{S}^{d}_{++}}
\newcommand{\PSD}{\mathbb{S}^{d}_{+}}
\newcommand{\normtwo}[1]{\left\lVert #1 \right\rVert_2}
\newcommand{\inner}[2]{\left\langle #1,#2 \right\rangle}
\newcommand{\rhoH}{\rho_{\mathrm{H}}}
\newcommand{\rhoW}{\rho_{\mathrm{W}}}

% \newtheorem{definition}{Definition}
% \newtheorem{remark}{Remark}
\usepackage{amsthm}
\newtheorem{remark}{Remark}

\usepackage{titlesec}

\titleformat{\chapter}[display]
  {\normalfont\huge\bfseries}
  {}{0pt}{\Huge}
\titlespacing*{\chapter}{0pt}{-40pt}{35pt}

% \newtcolorbox{questionbox}{
%   colback=gray!7,
%   colframe=black!45,
%   boxrule=0.5pt,
%   arc=2mm,
%   left=2mm,right=2mm,top=1.5mm,bottom=1.5mm
% }

% \newtcolorbox{resultbox}{
%   colback=blue!4,
%   colframe=blue!45!black,
%   boxrule=0.5pt,
%   arc=2mm,
%   left=2mm,right=2mm,top=1.5mm,bottom=1.5mm
% }

% -------------------------
% Figures, tables, floats
% -------------------------
\usepackage{graphicx}
\usepackage{subcaption}
\usepackage{booktabs}
\usepackage{array}
\usepackage{multirow}
\usepackage{longtable}
\usepackage{tabularx}
\usepackage{float}
\usepackage{placeins}
\usepackage{caption}
\captionsetup{font=small,labelfont=bf,justification=centering}
\graphicspath{{02_figures/}}

% A cleaner default for wide but not excessive mathematical plots
\newcommand{\thesisplot}[3][0.86\textwidth]{%
  \begin{figure}[tbp]
    \centering
    \includegraphics[width=#1]{#2}
    \caption{#3}
    \label{fig:#2}
  \end{figure}
}

% -------------------------
% Bibliography and references
% -------------------------
\usepackage[backend=biber,style=authoryear,maxcitenames=2,maxbibnames=99,giveninits=true,doi=true,url=true,isbn=false]{biblatex}
\usepackage{xcolor}
\definecolor{QMULBlue}{HTML}{003865}
\definecolor{ThesisGreen}{HTML}{006B54}
% Quiet thesis palette: pale fills, muted frames, and dark text.
\definecolor{ThesisInk}{HTML}{2D3E50}
\definecolor{TheoremTint}{HTML}{F0EEF7}
\definecolor{TheoremFrame}{HTML}{9A8BB5}
\definecolor{PropositionTint}{HTML}{EAF4F5}
\definecolor{PropositionFrame}{HTML}{78A7B0}
\definecolor{LemmaTint}{HTML}{EDF2F8}
\definecolor{LemmaFrame}{HTML}{8AA2BD}
\definecolor{DefinitionTint}{HTML}{EBF5F0}
\definecolor{DefinitionFrame}{HTML}{76A795}
\definecolor{AssumptionTint}{HTML}{FBF3E6}
\definecolor{AssumptionFrame}{HTML}{C5A06B}
\definecolor{ResultTint}{HTML}{F7EDF3}
\definecolor{ResultFrame}{HTML}{B58BA2}
\definecolor{ProofTint}{HTML}{F8F8F5}
\definecolor{ProofFrame}{HTML}{A5AAA5}
\definecolor{ProofTitle}{HTML}{E1E5E0}
\definecolor{ContributionTint}{HTML}{EAF5F0}
\definecolor{ContributionFrame}{HTML}{78B79A}
\definecolor{ContributionTitle}{HTML}{CBE5D8}
\usepackage[colorlinks=true,linkcolor=QMULBlue,citecolor=ThesisGreen,urlcolor=QMULBlue]{hyperref}
\usepackage[nameinlink,noabbrev,capitalise]{cleveref}

% Exact source helper for theorem titles, e.g.
% \begin{theorem}{Feynman--Kac formula \thmsource{\textcite[Theorem 8.2, p.~157]{oksendal}}}{fk}
\newcommand{\thmsource}[1]{\hfill\normalfont\footnotesize\emph{Source: #1}}

% -------------------------
% Theorems and boxed mathematical statements
% -------------------------
\usepackage[most]{tcolorbox}
\tcbuselibrary{theorems,breakable,skins}

\tcbset{
  thesisbox/.style={
    enhanced,
    breakable,
    colback=TheoremTint,
    colframe=TheoremFrame,
    coltitle=ThesisInk,
    fonttitle=\bfseries\fontsize{12}{14}\selectfont,
    attach boxed title to top left={xshift=1.5mm,yshift=-1.5mm},
    boxed title style={colback=TheoremTint,sharp corners,boxsep=0.6mm,left=0.6mm,right=0.6mm,top=0.2mm,bottom=0.2mm},
    sharp corners=south,
    boxrule=0.55pt,
    boxsep=0pt,
    left=1.2mm,right=1.2mm,top=6.2mm,bottom=1.4mm
  }
}

\newtcbtheorem[number within=chapter,crefname={theorem}{theorems},Crefname={Theorem}{Theorems}]{theorem}{Theorem}{thesisbox}{thm}
\newtcbtheorem[number within=chapter,crefname={proposition}{propositions},Crefname={Proposition}{Propositions}]{proposition}{Proposition}{thesisbox,colback=PropositionTint,colframe=PropositionFrame,boxed title style={colback=PropositionTint,sharp corners}}{prop}

\newtcbtheorem[number within=chapter,crefname={corollary}{corollary},Crefname={Corollary}{Corollary}]{corollary}{Corollary}{thesisbox,colback=white,colframe=black,boxed title style={colback=white,sharp corners}}{lem}

\newtcbtheorem[number within=chapter,crefname={lemma}{lemmas},Crefname={Lemma}{Lemmas}]{lemma}{Lemma}{thesisbox,colback=LemmaTint,colframe=LemmaFrame,boxed title style={colback=LemmaTint,sharp corners}}{lem}
\newtcbtheorem[number within=chapter,crefname={definition}{definitions},Crefname={Definition}{Definitions}]{definition}{Definition}{thesisbox,colback=DefinitionTint,colframe=DefinitionFrame,boxed title style={colback=DefinitionTint,sharp corners}}{def}
\newtcbtheorem[number within=chapter,crefname={assumption}{assumptions},Crefname={Assumption}{Assumptions}]{assumption}{Assumption}{thesisbox,colback=AssumptionTint,colframe=AssumptionFrame,boxed title style={colback=AssumptionTint,sharp corners}}{assump}
\newtcbtheorem[number within=chapter,crefname={result}{results},Crefname={Result}{Results}]{result}{Result}{thesisbox,colback=ResultTint,colframe=ResultFrame,boxed title style={colback=ResultTint,sharp corners}}{res}

% Boxed proof environment: use sparingly for important proofs.
% \newtcolorbox{proofbox}[1][]{
%   enhanced,breakable,colback=ProofTint,colbacktitle=ProofTitle,colframe=ProofFrame,coltitle=ThesisInk,title={Proof},fonttitle=\bfseries\fontsize{12.5}{15}\selectfont,boxrule=0.55pt,boxsep=0pt,left=1.2mm,right=1.2mm,top=1.8mm,bottom=1.4mm,lefttitle=1.2mm,righttitle=1.2mm,toptitle=1mm,bottomtitle=1mm,#1
% }

\newtcolorbox{proofbox}[1][]{
  enhanced,
  breakable,
  colback=ProofTint,
  colbacktitle=ProofTitle,
  colframe=ProofFrame,
  coltitle=ThesisInk,
  title={Proof},
  fonttitle=\bfseries\fontsize{12.5}{15}\selectfont,
  attach title to upper=false,   % <- forces the title to stay in its own bar,
                                  %    never gets prepended inline to the body text
  boxrule=0.55pt,
  boxsep=0pt,
  left=1.2mm,right=1.2mm,
  top=3mm,bottom=1.4mm,          % <- a touch more top padding so the first
                                  %    line of the proof clearly clears the bar
  lefttitle=1.2mm,righttitle=1.2mm,
  toptitle=1mm,bottomtitle=1mm,
  #1
}

% Chapter contribution summary box
\newtcolorbox{contributionbox}[1][Chapter contributions]{
  enhanced,breakable,colback=ContributionTint,colbacktitle=ContributionTitle,colframe=ContributionFrame,coltitle=ThesisInk,title={#1},fonttitle=\bfseries\fontsize{12.5}{15}\selectfont,boxrule=0.55pt,boxsep=0pt,left=1.2mm,right=1.2mm,top=1.8mm,bottom=1.4mm,lefttitle=1.2mm,righttitle=1.2mm,toptitle=1mm,bottomtitle=1mm
}

% Chapter epigraphs / italic quotes
\usepackage{epigraph}
\setlength{\epigraphwidth}{0.72\textwidth}
\renewcommand{\epigraphflush}{center}
\renewcommand{\textflush}{flushright}
\renewcommand{\epigraphsize}{\small\itshape}
\newcommand{\chapterquote}[2]{\epigraph{#1}{#2}}

% -------------------------
% Headings and page style
% -------------------------
\usepackage{titlesec}
% \titleformat{\chapter}[display]
%   {\normalfont\huge\bfseries}
%   {\chaptertitlename\ \thechapter}{18pt}{\Huge}
% \titlespacing*{\chapter}{0pt}{-5pt}{35pt}

\usepackage{fancyhdr}
\pagestyle{fancy}
\fancyhf{}
\fancyhead[LE]{\small\nouppercase{\leftmark}}
\fancyhead[RO]{\small\nouppercase{\rightmark}}
\fancyfoot[C]{\thepage}
\renewcommand{\headrulewidth}{0.4pt}

% -------------------------
% Utility commands
% -------------------------
\newcommand{\todo}[1]{\textcolor{red}{\textbf{TODO:} #1}}


\DeclareRobustCommand{\proofbox}[2][gray!10]{%
\begin{tcolorbox}[   %% Adjust the following parameters at will.
        breakable,
        left=0pt,
        right=0pt,
        top=0pt,
        bottom=0pt,
        colback=#1,
        %colframe=black,
        width=\dimexpr\textwidth\relax, 
        enlarge left by=0mm,
        boxsep=5pt,
        arc=0pt,outer arc=0pt,
        ]
        #2
\end{tcolorbox}
}
